\begin{answer}

\textbf{(b)(i) Benchmarking models A and B on GSM8K}

Both models were evaluated on all 100 problems in \texttt{gsm8k\_first\_100.jsonl}
using the provided \texttt{standard\_prompt\_template} and
\texttt{standard\_output\_extractor}. The metric is exact match: the extracted
numerical answer is compared against the \texttt{numerical\_answer} field, and a
generation from which no answer can be extracted is counted as incorrect.
Unparseable generations are also reported separately, since they measure
instruction-following rather than arithmetic ability.

\begin{center}
\begin{tabular}{lcccc}
\toprule
Model & Correct & Unparseable & Wrong answer & Accuracy \\
\midrule
A & 95 & 0 & 5  & \textbf{95\%} \\
B & 54 & 2 & 44 & \textbf{54\%} \\
\bottomrule
\end{tabular}
\end{center}

Model A scores \textbf{95\%} and model B scores \textbf{54\%}. The two models are
queried through the same interface with the same prompt, so the 41-point gap is
attributable entirely to their differing system prompts. Model B's errors are
overwhelmingly wrong answers (44) rather than malformed output (2), so it is
producing well-formatted but incorrect arithmetic rather than failing to follow
the output format.

\vspace{1em}
\textbf{(b)(ii) An improved prompt for model A}

\textbf{My prompt.}

\begin{lstlisting}[language={},numbers=none,basicstyle=\footnotesize\ttfamily,frame=single]
Solve the following problem using as many steps of reasoning as you need,
showing each calculation. Before you commit to an answer, check your work
carefully.

Problem: {question}

Now report your result. The LAST line of your response must be the final
numerical answer, formatted exactly like this:

#### <numerical_answer>

Rules for that last line:
- It must begin with the four characters #### followed by a space.
- Put only a plain number after "####" -- no LaTeX, no \boxed{}, no dollar
  signs, no units, no commas, no bold markers.
- Even if you already stated the answer during your reasoning, you must still
  repeat it on this final #### line. A bare number on its own is not
  acceptable.
- Write nothing at all after that line.

For example, if the answer is 42, the last line must read exactly:

#### 42
\end{lstlisting}

\textbf{Why I expected this to improve performance.}

The standard prompt contains the instruction ``with two or fewer steps of
reasoning.'' GSM8K problems are multi-step arithmetic word problems, so this is a
\emph{binding} constraint: it caps the number of intermediate results the model
can write down below what many of the problems actually require. Since a
decoder-only model does its arithmetic in the tokens it emits, suppressing those
tokens suppresses the computation itself. My prompt removes the cap, asks for
each calculation to be shown, and adds an explicit self-check step before the
model commits to an answer.

Lifting the reasoning cap, however, changes the register of the response from a
terse answer to a long worked solution --- and in that register the model's
strong habit is to close with LaTeX, e.g.\ \verb|$\boxed{42}$|, abandoning the
required \verb|####| marker entirely. An earlier draft of this prompt lost a
substantial fraction of its outputs to exactly this failure. The format rules in
the second half of the prompt exist to counter it, and they use three levers:
the format requirement is placed \emph{after} the problem rather than before it,
so it is the most recent instruction in context; it asks for the answer on the
\emph{last} line rather than the ``only'' line, which is compatible with showing
extensive reasoning instead of contradicting it; and it explicitly names the
competing formats (LaTeX, \verb|\boxed{}|, a bare number) that were observed to
displace the marker.

\textbf{Results.}

\begin{figure}[h]
\centering
\includegraphics[width=0.62\textwidth]{gsm8k_prompt_comparison.png}
\caption{Model A's accuracy on the 100 GSM8K problems under the standard prompt
and under my improved prompt.}
\end{figure}

\begin{center}
\begin{tabular}{lcccc}
\toprule
Prompt & Correct & Unparseable & Wrong answer & Accuracy \\
\midrule
Standard & 95 & 0 & 5 & 95\% \\
Improved & 98 & 0 & 2 & \textbf{98\%} \\
\bottomrule
\end{tabular}
\end{center}

\textbf{Did performance improve?}

Yes, but by a modest margin: accuracy rose from 95\% to \textbf{98\%}, a gain of
3 points. The gain is a genuine improvement in arithmetic rather than a
formatting artifact --- both prompts produced 0 unparseable generations in this
run, and the number of substantively wrong answers fell from 5 to 2. In other
words the improved prompt fixed 3 of the 5 problems the standard prompt got
wrong, and broke none of the 95 it got right.

Two caveats on the size of this effect. First, the margin is small relative to
sampling noise: at $p \approx 0.96$ with $n = 100$, the standard error of a
single run is about 2 points, so a 3-point difference measured once is
suggestive rather than conclusive. Establishing it properly would require
repeated runs with error bars, or a larger evaluation set. Second, and more
fundamentally, the baseline leaves very little room to improve --- with the
standard prompt already at 95\%, only 5 problems were available to fix, so the
maximum achievable gain was 5 points and this prompt captured 3 of them. A
ceiling effect of this kind means the experiment cannot strongly distinguish a
good prompt from a very good one; the remaining 2 failures may simply be
problems this model cannot solve at any prompt.

\end{answer}
